% GENERATED FILE. Edit canonical lessons, then run npm run generate:books.
% canonical-source-hash: fnv1a64:39a91ad10153d356
% canonical-lessons: TE-C255-myuziyam, TE-C255-vyayamasala, TE-C255-vedika, TE-C255-nagaram, TE-C255-raccabanda

\chapter{Museum, Gym, Stage, City, Village meeting platform}
\label{ch:te-museum-gym-stage-city-village-meeting-platform}
\hlchaptermodality{255}

\begin{chapteropening}
\textbf{By the end of this chapter:} \emph{I can say a museum, a gym, a stage, a city and the village meeting platform.}

Complete the last lesson of chapter 255: I can say a museum, a gym, a stage, a city and the village meeting platform.
\end{chapteropening}

\section[\texorpdfstring{\te{మ్యూజియం} (myūziyaṁ)}{మ్యూజియం (myūziyaṁ)}]{\te{మ్యూజియం} (myūziyaṁ) — a museum}

\label{lesson:TE-C255-myuziyam}



\begin{quote}
Before the new one: say the Telugu for east, then the Telugu for west.
\end{quote}

\subsection*{You'll want to know: \te{మ్యూజియం}}
\textbf{\te{మ్యూజియం}} — \emph{myūziyaṁ} — \textquotedblleft{}a museum\textquotedblright{}.

\begin{grammarlens}[title={What you've built}]
One more word for this chapter.
\end{grammarlens}

\subsection*{Your turn}
\begin{itemize}
\raggedright
  \item \emph{Say it:} \emph{myūziyaṁ}
  \item \emph{Say it:} \emph{myūziyaṁ}, once more
  \item \emph{Say it:} say \emph{myūziyaṁ}
  \item \emph{Recall:} say the Telugu for east, then the Telugu for west, then say \emph{myūziyaṁ} again
\end{itemize}

\begin{tcolorbox}[breakable,colback=teal!4,colframe=teal!35!black,title={Before you move on}]
What does \te{మ్యూజియం} mean? (\textquotedblleft{}A museum\textquotedblright{}.) Say it once more.
\end{tcolorbox}
\section[\texorpdfstring{\te{వ్యాయామశాల} (vyāyāmaśāla)}{వ్యాయామశాల (vyāyāmaśāla)}]{\te{వ్యాయామశాల} (vyāyāmaśāla) — a gym}

\label{lesson:TE-C255-vyayamasala}



\begin{quote}
Before the new one: say the Telugu for west, then the Telugu for a museum.
\end{quote}

\subsection*{You'll want to know: \te{వ్యాయామశాల}}
\textbf{\te{వ్యాయామశాల}} — \emph{vyāyāmaśāla} — \textquotedblleft{}a gym\textquotedblright{}.

\begin{grammarlens}[title={What you've built}]
One more word for this chapter.
\end{grammarlens}

\subsection*{Your turn}
\begin{itemize}
\raggedright
  \item \emph{Say it:} \emph{vyāyāmaśāla}
  \item \emph{Say it:} \emph{vyāyāmaśāla}, once more
  \item \emph{Say it:} say \emph{vyāyāmaśāla}
  \item \emph{Recall:} say the Telugu for west, then the Telugu for a museum, then say \emph{vyāyāmaśāla} again
\end{itemize}

\begin{tcolorbox}[breakable,colback=teal!4,colframe=teal!35!black,title={Before you move on}]
What does \te{వ్యాయామశాల} mean? (\textquotedblleft{}A gym\textquotedblright{}.) Say it once more.
\end{tcolorbox}
\section[\texorpdfstring{\te{వేదిక} (vēdika)}{వేదిక (vēdika)}]{\te{వేదిక} (vēdika) — a stage, a platform}

\label{lesson:TE-C255-vedika}



\begin{quote}
Before the new one: say the Telugu for a museum, then the Telugu for a gym.
\end{quote}

\subsection*{You'll want to know: \te{వేదిక}}
\textbf{\te{వేదిక}} — \emph{vēdika} — \textquotedblleft{}a stage, a platform\textquotedblright{}.

\begin{grammarlens}[title={What you've built}]
One more word for this chapter.
\end{grammarlens}

\subsection*{Your turn}
\begin{itemize}
\raggedright
  \item \emph{Say it:} \emph{vēdika}
  \item \emph{Say it:} \emph{vēdika}, once more
  \item \emph{Say it:} say \emph{vēdika}
  \item \emph{Recall:} say the Telugu for a museum, then the Telugu for a gym, then say \emph{vēdika} again
\end{itemize}

\begin{tcolorbox}[breakable,colback=teal!4,colframe=teal!35!black,title={Before you move on}]
What does \te{వేదిక} mean? (\textquotedblleft{}A stage, a platform\textquotedblright{}.) Say it once more.
\end{tcolorbox}
\section[\texorpdfstring{\te{నగరం} (nagaraṁ)}{నగరం (nagaraṁ)}]{\te{నగరం} (nagaraṁ) — a city}

\label{lesson:TE-C255-nagaram}



\begin{quote}
Before the new one: say the Telugu for a gym, then the Telugu for a stage.
\end{quote}

\subsection*{You'll want to know: \te{నగరం}}
\textbf{\te{నగరం}} — \emph{nagaraṁ} — \textquotedblleft{}a city\textquotedblright{}.

\begin{grammarlens}[title={What you've built}]
One more word for this chapter.
\end{grammarlens}

\subsection*{Your turn}
\begin{itemize}
\raggedright
  \item \emph{Say it:} \emph{nagaraṁ}
  \item \emph{Say it:} \emph{nagaraṁ}, once more
  \item \emph{Say it:} say \emph{nagaraṁ}
  \item \emph{Recall:} say the Telugu for a gym, then the Telugu for a stage, then say \emph{nagaraṁ} again
\end{itemize}

\begin{tcolorbox}[breakable,colback=teal!4,colframe=teal!35!black,title={Before you move on}]
What does \te{నగరం} mean? (\textquotedblleft{}A city\textquotedblright{}.) Say it once more.
\end{tcolorbox}
\section[\texorpdfstring{\te{రచ్చబండ} (raccabaṇḍa)}{రచ్చబండ (raccabaṇḍa)}]{\te{రచ్చబండ} (raccabaṇḍa) — the village meeting platform}

\label{lesson:TE-C255-raccabanda}



\begin{quote}
Before the new one: say the Telugu for a stage, then the Telugu for a city.
\end{quote}

\subsection*{You'll want to know: \te{రచ్చబండ}}
\textbf{\te{రచ్చబండ}} — \emph{raccabaṇḍa} — \textquotedblleft{}the village meeting platform\textquotedblright{}.

\begin{grammarlens}[title={What you've built}]
One more word for this chapter.
\end{grammarlens}

\subsection*{Your turn}
\begin{itemize}
\raggedright
  \item \emph{Say it:} \emph{raccabaṇḍa}
  \item \emph{Say it:} \emph{raccabaṇḍa}, once more
  \item \emph{Say it:} say \emph{raccabaṇḍa}
  \item \emph{Recall:} say the Telugu for a stage, then the Telugu for a city, then say \emph{raccabaṇḍa} again
\end{itemize}

\begin{tcolorbox}[breakable,colback=teal!4,colframe=teal!35!black,title={Before you move on}]
What does \te{రచ్చబండ} mean? (\textquotedblleft{}The village meeting platform\textquotedblright{}.) Say it once more.
\end{tcolorbox}
